% Geometría y márgenes
% paquetes para escribir codigo
\usepackage{listings}
\usepackage{xcolor}
% Márgenes y caja de texto
\usepackage[
  top=1.2cm,
  bottom=1.8cm,
  left=1.5cm,
  right=1.5cm,
  headheight=14pt,
  headsep=6pt,
  heightrounded
]{geometry}
\usepackage{booktabs}
\usepackage{bm}

% Idioma y codificaciones
%\usepackage[utf8]{inputenc}
%\usepackage[T1]{fontenc}
%\usepackage[spanish]{babel}

% Compilar con lualatex
\usepackage{fontspec} % no cargar inputenc ni fontenc con LuaLaTeX
\defaultfontfeatures{Ligatures=TeX}

% Fuentes con buen soporte Unicode (recomendado)
\setmainfont{Libertinus Serif}
\setsansfont{Libertinus Sans}
\setmonofont{Libertinus Mono}

% Cargar babel después de fontspec
\usepackage[spanish]{babel}

% Desactivar shorthands conflictivos de babel-es (comillas, <, >, .)
\AtBeginDocument{\shorthandoff{"<>.}}


\usepackage{tikz}
\usepackage{tcolorbox}
\usetikzlibrary{arrows.meta,angles,quotes,decorations.pathreplacing}
\usetikzlibrary{positioning, spy}
\usetikzlibrary{decorations.markings}



\usepackage{pgfplots}
\pgfplotsset{compat=1.18}
\usepgfplotslibrary{groupplots}


% Matemática y símbolos
\usepackage{amsmath,amssymb,amsfonts}
\usepackage{eucal}
\DeclareMathOperator{\tr}{Tr}

% Gráficos
\usepackage{graphicx}
\usepackage{epsfig}
\usepackage{pstricks}
\usepackage{pst-plot}

% Columnas y cajas
\usepackage{multicol}
\usepackage{fancybox}

% Estilo de página
%\pagestyle{empty}
%\headheight 0pt
%\textwidth 6.5in
%\textheight 9in
%\oddsidemargin 0in

% Utilidades
\usepackage{lipsum}
\usepackage{color}
\usepackage{hyperref}

% Comentarios y verbatim
\usepackage{verbatim} % verbatim environment

% -------------------------------------------------------
% Soluciones: ocultar/mostrar sin usar comment.cut
% -------------------------------------------------------
\usepackage{comment} % mantiene compatibilidad si lo necesitas en otras partes
\makeatletter
% Por defecto: ocultar soluciones (ni siquiera se parsean)
\excludecomment{solution}

% Si el Makefile define \WITHANSWERS, activamos un entorno 'solution'
% que imprime el contenido (sin pasar por comment.cut).
\@ifundefined{WITHANSWERS}{}{%
  \let\solution\relax
  \let\endsolution\relax
  \newenvironment{solution}{%
    \par\medskip
    \noindent\textbf{Solución.}\ %
  }{%
    \par\medskip
  }
  % Compatibilidad hacia atrás:
  % Si tus archivos viejos usan \begin{comment}...\end{comment} para soluciones,
  % mapeamos 'comment' -> 'solution' sólo cuando WITHANSWERS está activo.
  \let\comment\solution
  \let\endcomment\endsolution
}
\makeatother


% Listings (código)
\usepackage{listings}
\definecolor{codegreen}{rgb}{0,0.6,0}
\definecolor{codegray}{rgb}{0.5,0.5,0.5}
\definecolor{codepurple}{rgb}{0.58,0,0.82}
\definecolor{backcolour}{rgb}{0.95,0.95,0.92}
\lstdefinestyle{mystyle}{
  backgroundcolor=\color{backcolour},
  commentstyle=\color{codegreen},
  keywordstyle=\color{magenta},
  numberstyle=\tiny\color{codegray},
  stringstyle=\color{codepurple},
  basicstyle=\footnotesize\ttfamily,
  breakatwhitespace=false,
  breaklines=true,
  captionpos=b,
  keepspaces=true,
  numbers=left,
  numbersep=5pt,
  showspaces=false,
  showstringspaces=false,
  showtabs=false,
  tabsize=2
}
\lstset{style=mystyle}

\lstnewenvironment{pythonjupyter}{
  \lstset{
    language=Python,
    backgroundcolor=\color{backcolour},
    frame=lines,
    framerule=0.4pt,
    rulecolor=\color{codegray},
    basicstyle=\footnotesize\ttfamily,
    keywordstyle=\color{blue!70!black},
    commentstyle=\color{codegreen},
    stringstyle=\color{codepurple},
    numbers=left,
    numberstyle=\tiny\color{codegray},
    numbersep=6pt,
    breaklines=true,
    showstringspaces=false,
    tabsize=2
  }
}{}


% Macros del usuario (conservadas del .tex original)
\newcommand*{\op}[1]{\check{\mathbf#1}}
\newcommand{\bra}[1]{\langle #1 \rvert}
\newcommand{\ket}[1]{\lvert #1 \rangle}
\newcommand{\braket}[2]{\langle #1 \,\vert\, #2 \rangle}
\newcommand{\mean}[1]{\langle #1 \rangle}
\newcommand{\opvec}[1]{\check{\vec #1}}
\renewcommand{\sp}[1]{$${\begin{split}#1\end{split}}$$}
